%% kapa: exam=1 round=24 year=2013 period=2 form=A subject=accounting subject_ko=회계학 seq=11 number=11
%% source: Q-Net 감정평가사 기출문제 (cst00309) artlSeq=5073802 file=extracted/1008/5073802/2013년 제24회 감정평가사 1차시험 기출문제/2교시(A형).hwp
%% conversion: kordoc 4.24.51 -> tools/gen1.py (auto); review=proofread-sonnet; flags=none
\begin{question}{11}
(주)한국의 20×1년 상품과 관련된 내용은 다음과 같다. \par
{\small\centering
\begin{tabular}{|c|c|c|}
\hline
\makecell{구분} & \makecell{기초재고} & \makecell{당기매입} \\
\hline
\makecell{원가} & \makecell{￦14,000} & \makecell{￦51,000} \\
\hline
\makecell{판매가} & \makecell{15,000} & \makecell{85,000} \\
\hline
\end{tabular}\par}

20×1년도 매출(판매가)은 ￦74,000이고, 20×1년 말 상품의 순실현가능가치는 ￦16,000이다. (주)한국은 상품의 원가측정방법으로 소매재고법을 선택하였다. 원가흐름에 대한 가정으로 평균법을 적용하는 경우와 선입선출법을 적용하는 경우 각각의 평가방법에 따른 상품평가손실액의 차이는? (단, 평가손실충당금의 기초잔액은 없는 것으로 한다.) \par
\begin{choices}
  \item ￦400
  \item ￦600
  \item ￦900
  \item ￦1,000
  \item ￦1,300
\end{choices}
\end{question}
\answer{3}
